\section{Characteristic ideals over Krull domains}

The characteristic-ideal product formula for Krull domains is
classical~\cite[Proposition~2.1, author pp.~2--3]{Chabert2003Factorial}.
The criterion for a regular basis in terms of principal characteristic
ideals is classical as well~\cite[Proposition~2.4 and Corollary~2.5,
author p.~4]{Chabert2003Factorial}; see also
\cite[Proposition~5]{Chabert2014RegularBases}.
The proof below also computes the intersections with an arbitrary
coefficient overring. In the Noetherian normal case, the height-one
intersection and local projectivity inputs can be found in
\cite[Tags~031T and 00NX]{StacksNormalDomains}.

\begin{theorem}\label{thm:krull-characteristic-decomposition}
Let $R$ be a Krull domain with fraction field $K$, and let
$R\subset L\subset K$. For each height-one prime $\mathfrak p$ with
finite residue field, put $q_{\mathfrak p}=|R/\mathfrak p|$ and
\[
 w_q(n)=\sum_{j\ge1}\left\lfloor\frac n{q^j}\right\rfloor,
 \qquad I_n=\prod_{\substack{\operatorname{ht}\mathfrak p=1\\
                         |R/\mathfrak p|<\infty}}
                      \mathfrak p^{-w_{q_{\mathfrak p}}(n)}.
\]
Each product has finite support. There are monic polynomials
$P_n\in R[X]$, $\deg P_n=n$, such that
\[
 \operatorname{Int}(R)=\bigoplus_{n\ge0}I_nP_n,
 \qquad
 \Omega(R,L)=\bigoplus_{n\ge0}(I_n\cap L)P_n.
\]
Every $I_n\cap L$ is an invertible fractional ideal. Consequently
$\Omega(R,L)$ is projective and all its multivariable tensor maps
are isomorphisms. The same conclusion applies to every torsion-free
coefficient extension of an infinite Krull domain, after rational
coefficient reduction. In particular it applies to every Noetherian
normal domain.
\end{theorem}

\begin{proof}
The finite-field case is immediate, so assume that $R$ is infinite.
We use the defining height-one intersection and finite-character
properties of a Krull domain.

A height-one prime $\mathfrak p$ with finite quotient is maximal and
invertible. Choose $b\in R$ with $v_{\mathfrak p}(b)=1$. There are
only finitely many other height-one primes $\mathfrak q_i$ containing
$b$. Choose $a_i\in\mathfrak q_i\setminus\mathfrak p$ and put
$a=\prod_i a_i^{v_{\mathfrak q_i}(b)}$. Then $x=a/b$ has value $-1$
at $\mathfrak p$ and nonnegative value at every other height-one
prime. The intersection formula gives $x\mathfrak p\subset R$.
The ideal $\mathfrak p\mathfrak p^{-1}$ contains $\mathfrak p$ and
$xb=a\notin\mathfrak p$, so it is $R$. This proves invertibility
without a finite-generation assumption on arbitrary prime ideals.

For each fixed finite cardinality $q$, choose $a\in R$ with
$a^q-a\ne0$. Every prime with residue field of size $q$ contains
this nonzero element, so finite character gives only finitely many
height-one such primes. Since $w_q(n)=0$ for $q>n$, the product
defining $I_n$ is finite.

We recall the local coefficient calculation. In a DVR $V$ with
residue cardinality $q$, choose representatives $r_0=0,r_1,\ldots,r_{q-1}$
and a uniformizer $\pi$. For $n=\sum n_jq^j$ put
$a_n=\sum r_{n_j}\pi^j$ and $N_n(X)=\prod_{i<n}(X-a_i)$.
For every $x\in V$, counting the nodes in its congruence class at
each level gives
\[
 v(N_n(x))\ge\sum_{j\ge1}\lfloor n/q^j\rfloor=w_q(n),
 \qquad v(N_n(a_n))=w_q(n).
\]
Thus $N_n/\pi^{w_q(n)}$ is integer-valued, and its value at $a_n$
is a unit. Successive evaluation at $a_0,a_1,\ldots$ shows that
these normalized Newton polynomials form a regular $V$-basis.
In particular the degree-$n$ leading coefficients form
$\pi^{-w_q(n)}V$. The same argument works for polynomials evaluated
only on a subring containing the chosen nodes.

Apply this calculation at a finite-residue height-one prime of $R$.
Representatives and a valuation-one element can be chosen in $R$,
so all nodes belong to $R$. At every remaining height-one prime,
$R/\mathfrak p$ is infinite; a unit Vandermonde matrix on sufficiently
many $R$-points forces every coefficient to lie in $R_{\mathfrak p}$.
Hence the leading coefficient of an integer-valued polynomial of
degree at most $n$ belongs to $I_n$.

For the finitely many primes contributing to $I_n$, use the Chinese
remainder theorem coefficient by coefficient to choose one monic
$P_n\in R[X]$ of degree $n$ congruent to the corresponding $N_n$
modulo $\mathfrak p^{w_{q_{\mathfrak p}}(n)}R_{\mathfrak p}$.
The reduction map to this local quotient is surjective because
$\mathfrak p$ is maximal. It follows that $I_nP_n$ is integer-valued
at every height-one localization, and hence is integer-valued on $R$.
Subtracting the leading coefficient times $P_n$ and descending in
degree proves the first direct-sum decomposition. Since each $P_n$
is monic with coefficients in $R$, the same subtraction preserves
$L[X]$ and gives the second decomposition.

It remains to verify projectivity for $L$ without assuming that it is
finitely generated. At a contributing prime, the overring
$L_{\mathfrak p}$ of the DVR $R_{\mathfrak p}$ is either that DVR
or $K$. Localization commutes with the intersection of two submodules.
At every other maximal ideal, $(I_n)_{\mathfrak m}=R_{\mathfrak m}$.
Testing equality at all maximal ideals therefore gives
\[
 I_n\cap L=
 \prod_{\substack{\operatorname{ht}\mathfrak p=1,\ |R/\mathfrak p|<\infty\\
                              L_{\mathfrak p}=K}}
                \mathfrak p^{-w_{q_{\mathfrak p}}(n)}.
\]
This is a finite product of invertible fractional ideals. Their direct
sum is projective, and the finite dual-basis theorem gives all tensor
isomorphisms and plethory compatibility.
\end{proof}

\begin{corollary}\label{cor:regular-basis-principal-characteristic}
In the preceding theorem, $\Omega(R,L)$ has a regular basis, with
one basis polynomial in each degree, if and only if every ideal
$I_n\cap L$ is principal. In particular this holds when these ideals
are principal, for example over a local base in the stated class.
\end{corollary}
\begin{proof}
A generator $c_n$ gives the basis element $c_nP_n$. Conversely the
leading coefficient of the degree-$n$ member of a regular basis
generates the degree-$n$ characteristic ideal. This is a statement
about regular bases; it does not assert that an arbitrary free-module
basis must respect the degree filtration.
\end{proof}
